\section{How to Read This Manual}

This manual follows a calculation from its physical assumptions to its
numerical result. Each method begins with the physical picture and continuum
equations, then explains how the numerical representation gives a sequence
of updates. Initial and boundary data complete the problem. Source files,
diagnostics, and limitations connect that calculation to its implementation.
The later computing chapters explain files, restarts, parallel execution,
graphics-processor acceleration, plotting, and animation. Read the method
chapter first, then use the operating chapters to build and run it.

The first chapters describe the most exercised persistent-marker methods:
weighted particle-in-cell and spherical direct simulation Monte Carlo.  The
relativistic chapters then develop Einstein--Vlasov evolution in Schwarzschild
and ingoing Eddington--Finkelstein coordinates.  Vlasov--Poisson and
Boltzmann--Poisson follow.  Exact finite-angle collision physics is introduced
before small-angle approximations.  Semi-Lagrangian, discontinuous-Galerkin,
canonical-orbit, and other methods whose validation is still developing are
placed after the production modules.

\subsection{Physical and numerical conventions}

Collisionless relaxation redistributes orbits through the self-consistent
mean field while preserving the distribution along characteristics.
Collisional relaxation additionally redistributes momentum and energy through
the selected encounter law. Vlasov--Poisson, Boltzmann--Poisson and
Einstein--Vlasov denote the Newtonian collisionless, Newtonian collisional
and relativistic collisionless systems, respectively. Material
self-interacting dark matter (SIDM) and extended-subhalo encounters use
distinct cross sections and conservation budgets.

In a Newtonian module, \(\rho\) is represented mass density and
\(\rho_{\rm src}\) is the source density in Poisson's equation.
An isolated static halo uses total density. An expanding module uses its
stated density contrast or subtracts the prescribed homogeneous density
\(\bar\rho\). The expansion rate is \(H=\dot a_{\rm H}/a_{\rm H}\), where
\(a_{\rm H}\) is the cosmological scale factor and the dot is a physical-time
derivative. Physical radial velocity is \(u=Hr+w\), with peculiar radial
velocity \(w\) and radius \(r\). The gravitational potential is \(\Phi\).
Marker mass \(m_p\) and invariant rest weight \(W_p\) follow the convention
specified by the selected method.

Relativistic chapters specify the distribution normalization and momentum
measure on each time slice. Schwarzschild (SC) and ingoing
Eddington--Finkelstein (EF) coordinates use different momentum variables
and source measures. Their solved coordinate map identifies the same
physical events. A coordinate comparison therefore checks both the metric
and the transformed matter sources. The chapter notation distinguishes
rest mass, enclosed gravitational mass and the conserved outer
Arnowitt--Deser--Misner (ADM) mass.

Each method chapter defines its additional coefficients and integration
weights before their equations. The repository's \path{docs/TERMS.md}
collects the numerical and physical conventions used across the guides.

\subsection{Meaning of the maturity labels}

The labels below describe the calculations and checks available for each
method. A new physical problem requires its own convergence study, using the
observables and numerical controls relevant to that problem.

\begin{description}
\item[Production, strength-tested.] The task has a maintained executable,
      saved-state and diagnostic interface, representative successful-run
      material, and problem-specific checks.  Scientific use still requires
      resolution, timestep, domain, conservation, and stochastic-realization
      tests appropriate to that task.
\item[Validated reference.] The implementation is maintained as a comparison
      or baseline, but its coordinate patch, cost, dimensionality, or physical
      assumptions restrict the questions it can answer.
\item[Small-grid accuracy reference.] The method is intentionally expensive and is
      most useful on small grids for checking a faster approximation.
\item[Research or developing.] The algorithm is implemented and tested, but
      its production validation is incomplete or problem dependent.  A
      successful run is evidence of execution, not of mature predictive
      accuracy.
\item[Compatibility interface.] The entry point exists to interpret older
      inputs.  New work should use the task-specific executable named in the
      module catalogue.
\item[Not implemented.] The release has no such numerical path.  A similarly
      named option in another module must not be interpreted as support.
\end{description}

\begin{center}
\small
\begin{tabularx}{\linewidth}{@{}p{0.32\linewidth}p{0.25\linewidth}X@{}}
\toprule
Method family & Current label & What the label permits \\
\midrule
CPU Newtonian WPIC: collisionless, subhalo, and constant-cross-section
material SIDM tasks & Production, strength-tested within the evidence stated
in each module & Quantitative use still requires task-specific marker, mesh,
event-clock, locality, conservation, and realization tests. \\
Velocity-dependent material-SIDM WPIC & Research path & The collision law and
event scheduler are implemented, but the manual does not claim a completed
long-time gravothermal validation. \\
Constant-cross-section persistent-marker SphDSMC & Production benchmark path
with finite-time benchmark comparison & The documented comparison tests the stated
finite interval; it is not evidence for every cross-section law or for an
unresolved long-time core collapse. \\
Velocity-dependent persistent-marker SphDSMC & Research path & The source and
diagnostics are maintained, but a long-time core-validation claim is not made. \\
Schwarzschild Einstein--Vlasov WPIC & Validated exterior reference & Reliable
only while the polar coordinate gap remains positive; it does not evolve
through an apparent horizon. \\
Ingoing-EF Einstein--Vlasov WPIC & Research path with required
accuracy checks & Post-horizon interpretation requires the complete accepted-
step, provenance, causal-excision, and sensitivity evidence stated in its
chapter. \\
Nodal grid EV/VP and explicit finite-volume VP & Validated reference & Useful
for controlled grid studies within their documented transport and boundary
limits. \\
Dense exact finite-angle Boltzmann operator & Small-grid accuracy reference & Enumerates
the retained discrete collision channels; ``exact'' does not remove angular,
velocity-grid, remap, boundary, or time-discretization error. \\
Reversible, sampled, streamed, canonical-orbit, and SLDG methods & Research or
developing & Require the method-specific omission, projection, entropy,
sampling, carrier, and geometry checks before scientific interpretation. \\
Fokker--Planck/Landau operator & Maintained approximation & Represents a
small-angle approximation and must be checked against its locality and
coefficient-resolution assumptions. \\
GPU WPIC paths & Implemented acceleration with stated parity scope & Use only
for the tasks listed in the GPU chapter. Material-collision MPI runs
require checking how the result changes with the number of processes. \\
\bottomrule
\end{tabularx}
\end{center}

\subsection{Coordinates, phase space, and recurring symbols}

VENUS assumes spherical symmetry.  The areal radius \(r\) is defined so that
a sphere at that coordinate has area \(4\pi r^2\).  The coordinate \(t\)
denotes the time variable of the active module; the relativistic chapters use
\(\tau\) for Schwarzschild time and \(v\) for ingoing advanced time so that the
two cannot be confused.  In particular, \(\tau\) is not used for an orbital
phase in this manual.  The canonical-orbit chapter writes its normalized
radial phase as \(\zeta_{\rm orb}\in[0,1)\), while the Hamiltonian SLDG
chapter uses \(s_{\rm orb}\) for a local, time-like coordinate along an
orbit.  Angles on the unit sphere are \(\theta\) and \(\phi\), and
\(\dd\Omega^2=\dd\theta^2+\sin^2\theta\,\dd\phi^2\).

The distribution function \(f\) is phase space density.  Its arguments depend
on the module, but \(w\) normally denotes radial momentum or peculiar radial
velocity, \(L=|\bm r\times\bm v|\) is angular momentum magnitude, and
\(\eta=L^2/2\) is the angular coordinate stored by many solvers.  A marker
index is \(p\); radial, radial-momentum, and angular-grid indices are
\(i,j,k\); and \(n\) labels a time level.  A half index such as \(i+1/2\)
denotes a cell face.  Every chapter restates the precise physical meaning and
units of these quantities before using its equations.

Unless a chapter says otherwise, VENUS uses geometrized relativistic units
\(G=c=1\), or the explicit Newtonian gravitational constant \(G\) in a
Newtonian module.  Here \(c\) is the speed of light.  The symbol \(m(r)\)
denotes enclosed gravitational mass. The symbols \(m_p\) or \(W_p\) denote
the positive mass or invariant rest-number weight carried by marker \(p\),
with the convention specified in its module. Enclosed gravitational mass
includes the energy and binding contributions appropriate to the field
equations; marker weights specify the represented matter.

\subsection{Indices, derivatives, and operators}

An ordinary derivative along a characteristic is written \(\dot x=\dd x/\dd t\) or with the time coordinate
named explicitly.  A prime denotes the radial derivative only where that
convention is stated.  The symbol \(\nabla_a\) is the spacetime covariant
derivative: it differentiates tensors while accounting for the changing
coordinate basis.  Lower-case Latin spacetime indices such as \(a,b\) run
over four coordinates; repeated tensor indices are summed.  Module-specific
two-dimensional indices are defined where introduced.

The collision operator is written \(\mathcal C[f]\): it is the rate at which
collisions change \(f\).  The symbol \(\Delta\) denotes a finite grid width or
finite change. A norm measures the size of a quantity, and a residual measures
how closely a numerical result satisfies a stated relation. Each diagnostic
chapter gives its normalization and acceptance threshold. The threshold sets
the required numerical accuracy for that calculation.

\subsection{Acronyms and numerical vocabulary}

\begin{description}
\item[EV, VP, and BP] Einstein--Vlasov, Vlasov--Poisson, and
      Boltzmann--Poisson.
\item[SC and EF] Schwarzschild coordinates and ingoing
      Eddington--Finkelstein coordinates.  In the Schwarzschild description,
      the spherical metric is diagonal and a sphere at radius \(r\) has
      area \(4\pi r^2\).  The metric functions evolve with the matter.
\item[PIC and WPIC] particle-in-cell and weighted particle-in-cell.  Markers
      move continuously, while their positive represented weights are
      deposited on a field mesh.
\item[CIC and TSC] cloud-in-cell and triangular-shaped-cloud deposition.  They
      are respectively two-point linear and three-point quadratic marker
      shapes on the selected mesh coordinate.
\item[DSMC and SphDSMC] direct simulation Monte Carlo and its persistent-
      marker spherical implementation.  ``Monte Carlo'' means that controlled
      random sampling estimates collision events.
\item[CDM and SIDM] collisionless cold dark matter and self-interacting dark
      matter.
\item[FP] Fokker--Planck, the small-angle drift--diffusion approximation to a
      collision operator.
\item[SLDG and DG] semi-Lagrangian discontinuous Galerkin and discontinuous
      Galerkin.  A discontinuous Galerkin method represents the solution by a
      polynomial in each cell without forcing adjacent cell polynomials to
      share the same boundary value.
\item[CFL] Courant--Friedrichs--Lewy safety factor, which limits a timestep by
      the fastest resolved motion or source timescale.
\item[RK4, KDK, RKL2, and NTC] classical fourth-order Runge--Kutta;
      kick--drift--kick; second-order Runge--Kutta--Legendre super-time-
      stepping; and no-time-counter collision sampling.
\item[RK, AB, SSPRK2, and STS] Runge--Kutta; Adams--Bashforth;
      two-stage strong-stability-preserving Runge--Kutta; and super-time-
      stepping.  These are time-integration methods, not new physical models.
\item[COM, RNG, MFP, and CDF] center-of-mass frame; random-number generator;
      mean free path; and cumulative distribution function.
\item[RHS, RMS, and DF] right-hand side of an equation; root-mean-square;
      and distribution function.
\item[LC] Levi--Civita regularization, used to retry a numerically stiff close
      passage in a transformed radial coordinate.
\item[ADM and GR] Arnowitt--Deser--Misner mass measured at spatial infinity,
      and general relativity.
\item[FLRW and TOV] Friedmann--Lema\^itre--Robertson--Walker cosmological
      geometry, and the Tolman--Oppenheimer--Volkoff stellar-structure
      equations.
\item[IC and I/O] initial condition, and input/output.
\item[CUDA, OpenMP, and MPI] NVIDIA's GPU programming platform; the OpenMP
      shared-memory model; and the Message Passing Interface used by
      distributed solver tasks.
\item[API and CLI] application programming interface and command-line
      interface.
\item[CSV, JSON, SHA-256, PNG, and GIF] comma-separated-value table;
      JavaScript Object Notation record; the SHA-256 cryptographic content
      fingerprint; Portable Network Graphics image; and Graphics Interchange
      Format animation.
\item[ODE and PDE] ordinary and partial differential equation.
\item[NFW and EdS] Navarro--Frenk--White halo profile and Einstein--de Sitter
      cosmological background.
\item[DAG and RAM] directed acyclic graph, a dependency graph with no closed
      loop; and random-access memory, the working memory available to a CPU or
      GPU process.
\item[IEEE 754] the standard representation and arithmetic rules used for the
      floating-point numbers stored by the supported binaries.
\item[IAS and IIT] Institute for Advanced Study, and Indian Institute of
      Technology.  Their names identify documented computing sites; this
      manual does not assume that a particular queue or software module is
      currently available.
\item[INI and QA] the conventional name for a plain key--value configuration
      file, and quality assurance.
\end{description}

\subsection{Plain-language numerical vocabulary}

\begin{description}
\item[Characteristic.] A trajectory in phase space along which the
      collisionless distribution is carried.  It is the path of one ideal
      phase space point under the current gravitational field.
\item[Grid, mesh, node, cell, and face.] A grid or mesh is a finite set of
      locations used to represent a continuous problem.  A node is a stored
      point, a cell is the interval or volume between neighboring faces, and a
      face is a cell boundary.  A nodal value and a cell average are different
      numerical objects even when they use the same coordinates.
\item[Deposit and interpolate.] Deposition spreads marker-carried matter to
      nearby mesh locations.  Interpolation performs the complementary task
      of estimating a field or distribution between stored locations.  The
      chosen weights determine conservation, smoothness, and boundary behavior.
\item[Quadrature.] A weighted finite sum that approximates an integral.  Its
      nodes and weights are part of the numerical resolution.
\item[Finite volume and flux.] A finite-volume method evolves the amount
      stored in each cell.  A flux is the amount crossing a cell face per unit
      time; neighboring cells exchange the same face flux with opposite signs.
\item[Remap, reconstruction, pullback, and stencil.] A reconstruction builds a
      within-cell approximation from stored data.  A remap transfers that
      approximation to another grid or coordinate.  A pullback evaluates the
      earlier solution at a traced departure point.  A stencil is the local set
      of stored values used by one such calculation.
\item[Majorant and accept--reject sampling.] A majorant is a proved upper
      bound on an event rate.  A trial drawn from that upper rate is accepted
      with the ratio of the true rate to the bound.  The upper bound guarantees a valid acceptance probability. A violation
      therefore stops the sampler and requires the rate bound to be corrected.
\item[Kernel, backend, and frontend.] A kernel is a focused numerical
      operation.  A backend is one selectable implementation of an operation.
      A frontend is the task-specific executable and parser that validates a
      request before calling shared kernels.
\item[Cadence.] The time or step interval between repeated events such as
      field reconstruction, collision application, output, or restart writing.
\item[Resolution and convergence.] Resolution states how finely a continuum
      is represented.  A convergence study repeats the calculation with finer
      representations or smaller timesteps and measures whether the answer
      approaches a stable limit.
\item[Locality.] A collision approximation is local only when paired objects
      occupy a region small compared with the physical variation or mean-free-
      path scale required by that method.
\item[Atomic write.] Data are first written to a temporary object and made
      visible under the final name only after the complete write succeeds.
      This prevents a partial final file from looking complete.
\end{description}

\subsection{Data objects and evidence words}

An \emph{initial-condition file} contains the state from which a fresh run
starts.  A \emph{checkpoint} or \emph{restart} contains enough evolved state
to continue a run.  A \emph{profile} is a saved radial field table; a
\emph{history} is a time-indexed diagnostic table; and a \emph{ledger} records
conserved or transferred quantities before and after an operation.  A
\emph{summary} is the machine-readable final status.  A \emph{manifest}
records identity, parameters, relationships, and content fingerprints for one
or more files.  A \emph{sidecar} is a companion file whose state is checked
against the primary checkpoint or output.  A \emph{certificate} is a
fail-closed record assembled only after its stated evidence gates pass; it is
not created from a visual judgement.

``Fail closed'' means that missing, stale, inconsistent, or unauthenticated
evidence produces failure or ``inconclusive,'' never a pass.  ``Provenance''
means the recorded chain from source, executable, parameters, initial data,
and parent artifacts to a result.  ``Deterministic'' means that the stated
inputs and execution contract reproduce the same discrete choices; it does
not mean that truncation or sampling error is absent.

\subsection{How code listings are identified}

Every displayed code or command listing is preceded by a \textbf{Listing
ownership} line.  Its \textbf{Task or tool} field names the executable,
library consumer, or analysis utility to which the listing belongs.  Its
\textbf{Implementation} field gives the repository path that owns the shown
behavior.  Its \textbf{Role} field says whether the listing is an invocation,
configuration fragment, data schema, or shortened implementation excerpt.
A shortened excerpt explains logic; the source file, not the excerpt, remains
authoritative.

Equation prefixes identify the numerical method: for example, EWP denotes
Einstein--Vlasov weighted particles and IEF denotes ingoing
Eddington--Finkelstein evolution. Section titles and cross-references connect
the governing equations to their discretization and implementation.
